\begin{spacing}{1}
    \chapter*{Abstract}
\end{spacing}

This thesis presents and documents new features and improvements to the internal testing system of the KEBA Group AG, an Austrian technology company specializing in automation systems for industry, banking and energy applications. 

After analyzing the initial problems and how to solve them, various used technologies and concepts are explained. This includes generally known systems such as Python and PostgreSQL, as well as more specific concepts such as the Robot Framework and the Monaco Editor.

Next, the thought process and implementation of the new features and improvements is described in detail. The new functionalities include the creation of a new EBNF-based test script format, the definiton of a new set of keywords for the test scripts and the development of a new alternative database system for the test data. 



\newpage
\begin{spacing}{1}
    \chapter*{Zusammenfassung}
\end{spacing}

Diese Diplomarbeit präsentiert und dokumentiert neue Funktionen und Verbesserungen am internen Testsystem der KEBA Group AG, einem österreichischen Technologieunternehmen, das sich auf Automatisierungssysteme für Industrie, Banken und Energieanwendungen spezialisiert.

Nach der Analyse der anfänglichen Probleme und deren möglichen Lösungen werden verschiedene verwendete Technologien und Konzepte erklärt. Dazu gehören allgemein bekannte Systeme wie Python und PostgreSQL sowie spezifischere Konzepte wie das Robot Framework und der Monaco Editor.

Anschließend wird der Denkprozess und die Implementierung der neuen Funktionen und Verbesserungen im Detail beschrieben. Zu den neuen Funktionalitäten gehören die Erstellung eines neuen EBNF-basierten Testskriptformats, die Definition von neuen Schlüsselwörtern für die Testskripte und die Entwicklung eines neuen alternativen Datenbanksystems für die Testdaten.
